\scanpage{213}

\mspage{87}

somme amalgamée $C=A_1\amalg_{A_0} A_2$, et le morphisme canonique $C\longrightarrow B$ est \deleted{un} mono.] \inserted{Supposons de plus que $M$ satisfait à la condition \textcircled{e} p.~76, i.e. \unsure{une|que} \deleted{les} $\varinjlim$ filtr. de monos est mono.} \deleted{Alors \unsure{on a ce qui suit}}

$1^\circ)$ Soit $\mathrm{Mon}\subset\mathrm{Fl}\,M$ l'ens des monomorphismes, \deleted{\unsure{alors}}, et $\mathrm{Mon}_0$ le sous-ensemble de $\mathrm{Mon}$ formé des \deleted{\unsure{mono}} $i\in\mathrm{Mon}$ tels que source et but appartiennent à \deleted{\unsure{$M$}} $M_0$. On a alors \deleted{\unsure{ce}}

% slanted
\margin{\uline{NB} On a là un critère d'injectivité d'objets \deleted{\unsure{dans}} des catégories $M_{/Y}$ -- il suffit de tester la condition d'\unsure{inj.} par des $i$ mono. \unsure{tels que|i.e. que} $i$ soit dans $\mathrm{Mon}_0$.}

\[
\mathrm{Mon}\deleted{\unsure{=}}\subset\widetilde{\mathrm{Mon}_0},\quad \text{d'où }\left\{\begin{aligned}
\widetilde{\mathrm{Mon}}&=\widetilde{\mathrm{Mon}_0}\\
(\mathrm{Mon})_*&=(\mathrm{Mon}_0)_*\ \deleted{\unsure{(*)}}
\end{aligned}\right.\tag*{(7.1.2)}
\]

\deleted{et si il satisfait à la condition \textcircled{e} page 76 \unsure{lim.} \unsure{et de la lim} \unsure{filtr.} \unsure{de i mono est mono})}

\inserted{$\mathrm{Mon}$ est stable \deleted{\unsure{par}} cochangements de base, et \uline{$M_0$ est petite}}, et si \deleted{\illegible{1} \unsure{mêmes}}

\deleted{\[ \mathrm{Mon}=\widetilde{\mathrm{Mon}_0}\tag*{(7.1.3)}
\]}

% margin note written sideways
\margin{si $\mathrm{Mon}$ est stable par co-changement de base}

$2^\circ)$ Si \inserted{$M_0$ est \uline{petit} et si} les objets de $M_0$ sont accessibles \inserted{(p. ex. $M$ satisfait \textcircled{c} p.~76)}, et \deleted{si $M$ satisfait \textcircled{\unsure{e|c}} ($\mathrm{Mon}$ stable par \uline{$\lim$} filtrantes)}, alors il existe une factorisation fonctorielle des \deleted{tout} flèches $f$ de $M$ en

\[
f=pi,
\]

avec $i\in\mathrm{Mon}$, $p\in(\mathrm{Mon})_*$. (\uline{NB} Une flèche $p\colon X\longrightarrow Y$ de $M$ est dans $(\mathrm{Mon})_*$, sss elle fait de $X$ un objet \uline{injectif} de $M_{/Y}$ -- donc l'assertion de factorisation assure que dans les catégories induites $M_{/Y}$ ``il y a assez d'injectifs''). De plus, \deleted{\unsure{on}} on a

\[
\mathrm{Mon}=\widetilde{\mathrm{Mon}_0}\tag*{(7.1.3)}
\]

$3^\circ)$ Soit $W$ une partie de $\mathrm{Fl}(M)$ qui est \uline{$L$-accessible} \unsure{lire|i.e.} accessible, et stable par limites inductives filtrantes. [On suppose à présent que $M$ satisfait les conditions a) b) c) d) e) de p.~76, d'où une bonne théorie des parties $L$-accessibles de $\operatorname{Ob}M$ ou de $\mathrm{Fl}\,M$.] On suppose, \inserted{plus précisément,} \deleted{de plus} que $W$ satisfait les conditions
% archive sheet number: 212
